%
% phasensprung.tex -- 
%
% (c) 2021 Prof Dr Andreas Müller, OST Ostschweizer Fachhochschule
%
\documentclass[tikz]{standalone}
\usepackage{amsmath}
\usepackage{times}
\usepackage{txfonts}
\usepackage{pgfplots}
\usepackage{csvsimple}
\usetikzlibrary{arrows,intersections,math,calc}
\definecolor{darkred}{rgb}{0.8,0,0}
\definecolor{darkgreen}{rgb}{0,0.6,0}
\begin{document}
\def\skala{1.2}
\begin{tikzpicture}[>=latex,thick,scale=\skala]

	% Achsen
	\draw[->] (-0.8,0) -- (5.4,0) node[right] {$\operatorname{Re}(z)$};
	\draw[->] (0,-1.4) -- (0,1.8) node[above] {$\operatorname{Im}(z)$};
	\draw[very thick,red] (0,0) -- (5.0,0);
	% Parameter
	\def\eps{0.6}
	
	% oberer und unterer Teil der Hankel-Kontur
	\draw[->,thick] (\eps,0.12) -- (4.8,0.12);
	\draw[->,thick] (4.8,-0.12) -- (\eps,-0.12);
	
	\node[above] at (3.2,0.18) {$x^{s-1}e^{-i\pi(s-1)}$};
	\node[below] at (3.2,-0.18) {$x^{s-1}e^{+i\pi(s-1)}$};

\end{tikzpicture}
\end{document}

